\chapter{Erythrasma}
\index{Erythrasma}

\section*{Definition and Overview}
Erythrasma is a common, harmless skin infection caused by the bacterium Corynebacterium minutissimum. It produces well-defined, reddish-brown patches of skin, most often in the armpits, groin, and between the toes, where skin rubs together and stays warm and moist. Erythrasma is more common in adults, in people with diabetes, and in warm, humid climates. It is easily treated with antibiotics, either as a cream or a tablet, and it is not contagious in the usual sense.

\section*{Epidemiology}
Erythrasma is a common condition, particularly in tropical and subtropical regions and in people who sweat heavily or are overweight. It affects adults more than children and is more common in men. The infection is frequently mistaken for a fungal infection such as tinea, and it may coexist with other skin conditions. It is easily treated, but it can recur, especially when the underlying conditions that favour it remain.

\section*{Aetiology and Pathophysiology}
Erythrasma is caused by a bacterium that thrives in warm, moist skin folds.
\begin{itemize}
    \item \textbf{The organism:} Corynebacterium minutissimum, a bacterium that normally lives on the skin
    \item \textbf{Conditions that favour it:} warmth, moisture, friction, sweating, obesity, and diabetes
    \item \textbf{Mechanism:} the bacteria multiply in the skin folds and produce a pigment, causing the characteristic red-brown patches
    \item \textbf{Location:} the armpits, groin, between the toes, and under the breasts
    \item \textbf{Distinguishing feature:} under a special UV light (Wood's lamp), the patches glow coral red
    \item \textbf{Not contagious:} it does not spread easily from person to person
\end{itemize}

\section*{Clinical Features}
Erythrasma has a distinctive appearance.
\begin{itemize}
    \item Well-defined, flat, reddish-brown or pink patches
    \item Fine scaling of the patches
    \item Itching, which is usually mild
    \item Common sites: armpits, groin, between the toes, and under the breasts
    \item In the toe webs, it may look like soggy, macerated skin
    \item A coral-red glow under a Wood's lamp
    \item Asymptomatic in many people
\end{itemize}

\section*{Differential Diagnosis}
Other skin conditions look similar.
\begin{itemize}
    \item Tinea cruris and tinea corporis, fungal infections that are more scaly and itchy
    \item Candidiasis, a yeast infection with satellite spots
    \item Intertrigo, simple chafing and inflammation of the skin folds
    \item Seborrhoeic dermatitis
    \item Psoriasis of the skin folds
    \item Contact dermatitis
\end{itemize}

\section*{When to Seek Medical Care}
Persistent, itchy, or spreading patches in the skin folds should be assessed by a GP. Diagnosis is usually straightforward and treatment effective.

\textbf{Contact your local emergency services for:}
\begin{itemize}
    \item Signs of a spreading bacterial skin infection, including fever and a rapidly enlarging red area
    \item Painful, weeping, or blistered skin
    \item Signs of sepsis in a person with diabetes
    \item A very unwell person
\end{itemize}

\section*{Diagnosis and Investigations}
Diagnosis is usually clinical.
\begin{itemize}
    \item \textbf{Examination:} the appearance and location of the patches
    \item \textbf{Wood's lamp:} a special UV light that makes the patches glow coral red, confirming the diagnosis
    \item \textbf{Skin scraping:} examined under the microscope to distinguish erythrasma from fungus
    \item \textbf{Blood tests:} for diabetes if this is suspected
\end{itemize}

\section*{Management}
Treatment is simple and effective.
\begin{itemize}
    \item \textbf{Antibiotic cream:} applied to the affected skin for a couple of weeks
    \item \textbf{Oral antibiotics:} for extensive or recurrent infection
    \item \textbf{Antibacterial wash:} helps in recurrent cases
    \item \textbf{Keeping the skin dry:} reducing moisture is important for cure and prevention
    \item \textbf{Treating the cause:} managing diabetes and weight where relevant
    \item \textbf{Follow-up:} most cases clear fully, but recurrence is possible
\end{itemize}

\section*{Complications}
\begin{itemize}
    \item Recurrence, especially in warm, moist conditions
    \item Coinfection with fungi, requiring combined treatment
    \item Chronic infection in people with diabetes or obesity
    \item Misdiagnosis as tinea, leading to ineffective antifungal treatment
\end{itemize}

\section*{Prognosis and Outcomes}
Erythrasma responds well to treatment, and most cases clear within a few weeks with antibiotic therapy. The main issue is recurrence in people who remain prone to it, which is managed by keeping the skin dry, treating any underlying condition, and retreating early if it returns.

\section*{Prevention and Self-Care}
\begin{itemize}
    \item Keep skin folds clean and dry
    \item Wear loose, breathable clothing
    \item Use antiperspirants to reduce sweating
    \item Dry thoroughly after showering and swimming
    \item Manage weight and diabetes
    \item Avoid sharing towels
    \item Seek early treatment for recurrence
\end{itemize}

\section*{Sources and Further Reading}
The following reputable sources provide further information for healthcare students and patients seeking reliable, current advice about erythrasma.
\begin{itemize}
    \item DermNet. \textit{Erythrasma.} https://dermnetnz.org/
    \item Healthdirect Australia. \textit{Skin conditions information.} https://www.healthdirect.gov.au/
    \item Australasian College of Dermatologists. \textit{Skin condition information.} https://www.dermcoll.edu.au/
    \item NHS. \textit{Erythrasma.} https://www.nhs.uk/conditions/erythrasma/
    \item Mayo Clinic. \textit{Erythrasma: overview.} https://www.mayoclinic.org/diseases-conditions/
    \item MSD Manual. \textit{Erythrasma.} https://www.msdmanuals.com/
\end{itemize}
